\begin{theorem}[\texorpdfstring{$q$}{q}-重同步直积的置换对称与碰撞核的内禀对易子]\label{thm:pom-momq-permutation-symmetry}
沿用定理 \ref{thm:pom-fiber-multiplicity-matrix-product} 与命题 \ref{prop:pom-momq-tensor-kernel-construction} 的记号。
令 \(V:=\CC^{\mathcal S}\)，并令 \(T_b:V\to V\)（\(b\in\{0,1\}\)）为相应的线性算子；对整数 \(q\ge 2\) 令 \(W:=V^{\otimes q}\)。
定义置换表示
\[
\pi:S_q\to \mathrm{GL}(W),
\qquad
\pi(\sigma)(v_1\otimes\cdots\otimes v_q)
=v_{\sigma^{-1}(1)}\otimes\cdots\otimes v_{\sigma^{-1}(q)}.
\]
则对每个 \(b\in\{0,1\}\) 与 \(\sigma\in S_q\) 成立交换关系
\[
\pi(\sigma)\,T_b^{\otimes q}=T_b^{\otimes q}\,\pi(\sigma).
\]
从而
\[
I_2\otimes \pi(S_q)\ \subseteq\ \mathrm{Comm}\bigl(A_q^{\mathrm{ten}}\bigr),
\]
即置换对称作为 \(q\)-重同步直积的结构性对易子，内禀地进入张量分解型碰撞核 \(A_q^{\mathrm{ten}}=G_0\otimes T_0^{\otimes q}+G_1\otimes T_1^{\otimes q}\)。
\end{theorem}
\begin{proof}
对任意 \(v_1,\dots,v_q\in V\)，由张量积函子性
\[
T_b^{\otimes q}(v_1\otimes\cdots\otimes v_q)=T_bv_1\otimes\cdots\otimes T_bv_q
\]
可见：先对张量因子作置换、再施加 \(T_b^{\otimes q}\) 与先施加 \(T_b^{\otimes q}\)、再置换完全一致，故 \(\pi(\sigma)\) 与 \(T_b^{\otimes q}\) 对易。
于是 \(I_2\otimes \pi(\sigma)\) 分别与 \(G_0\otimes T_0^{\otimes q}\) 与 \(G_1\otimes T_1^{\otimes q}\) 对易，从而与其和 \(A_q^{\mathrm{ten}}\) 对易。
\end{proof}
\leanverified{paper\_pom\_momq\_permutation\_symmetry}

\begin{theorem}[Schur--Weyl 最大性：不可约转移代数下的置换对易子刻画]\label{thm:pom-momq-schur-weyl-commutant}
沿用上述记号，并令
\[
\mathcal A:=\mathrm{Alg}_{\CC}\langle T_0,T_1\rangle\ \subseteq\ \mathrm{End}(V)
\]
为由逆码转移算子生成的复代数。
若 \(\mathcal A\) 在 \(V\) 上作用不可约，则对每个整数 \(q\ge 2\) 有
\[
\mathrm{Comm}_{\mathrm{End}(V^{\otimes q})}\!\bigl(\mathcal A^{\otimes q}\bigr)
\;=\;
\pi\bigl(\CC[S_q]\bigr),
\]
亦即：在张量因子上与 \(T_0^{\otimes q},T_1^{\otimes q}\)（等价地与 \(\mathcal A^{\otimes q}\)）同时对易的全部算子恰由置换作用生成。
\end{theorem}
\begin{proof}
由 Burnside 定理，\(\mathcal A\) 在 \(V\) 上不可约推出 \(\mathcal A=\mathrm{End}(V)\)。
于是 \(\mathcal A^{\otimes q}=\mathrm{End}(V)^{\otimes q}\) 在 \(V^{\otimes q}\) 上的对易代数由 Schur--Weyl 对偶性给出，正是置换表示的像 \(\pi(\CC[S_q])\)。
\end{proof}
\leanverified{paper\_pom\_momq\_schur\_weyl\_commutant}
